% =========================================================================
% Open-World RF Emitter Discovery — capstone/internship paper DRAFT v0.1
% Sources of every number (verbatim, no re-derivation):
%   summer_work/EVAL_PROTOCOL.md          (locked protocol, WiSig/DRFF cells, TEST)
%   gen_rff/FINDINGS.md                   (F1-F8, DANN NULL, generalist closed)
%   gen_rff/ext_protocols/EXT_FINDINGS.md (EF1-EF8, M1/M2, W1-W5, BLE D2)
%   gen_rff/demo_router/README.md         (router demo, S1-S6, gates)
%   summer_work/results/demo_play1b/DEMO_OPERATING_POINT.md
% Compile: IEEEtran (Overleaf / TeX Live). Not compiled locally (no pdflatex).
% =========================================================================
\documentclass[conference]{IEEEtran}
\usepackage{amsmath,amssymb}
\usepackage{booktabs}
\usepackage{graphicx}
\usepackage{multirow}
\usepackage{xcolor}
\usepackage[hidelinks]{hyperref}

\newcommand{\ari}{\mathrm{ARI}}
\newcommand{\todo}[1]{\textcolor{red}{[TODO: #1]}}

\begin{document}

\title{Open-World Discovery of Unseen RF Emitters:\\
Burst Integration, Native Enrollment, and the Limits of Generalist Fingerprint Encoders}

\author{\IEEEauthorblockN{Abhinav W}
\IEEEauthorblockA{\todo{affiliation, advisor, email}}}

\maketitle

\begin{abstract}
Radio-frequency (RF) fingerprinting research overwhelmingly evaluates \emph{rejection}: deciding
whether a signal comes from an enrolled transmitter. We study the harder open-world task of
\emph{discovery}: given emissions from devices never seen in training, group them into distinct
device identities and estimate how many there are. Across four device populations spanning three
protocols and sample rates from 5 to 50\,MS/s --- 150 WiFi transmitters (WiSig), 16 USRP X310s
(ORACLE), 8 same-model DJI airframes (DRFF), and 31 same-model BLE units --- we report three main
results under a pre-registered, device-disjoint, leakage-audited protocol. \textbf{(1) The
cross-domain transfer wall is integration-dependent, not fixed:} at a burst size of $N{=}120$
windows, a frozen WiFi encoder, 19 hand-crafted physics features, and a natively trained encoder
all converge into a single 0.62--0.79 ARI band on drone discovery; the residual advantage of
native training is only $\approx$0.06--0.11 ARI. \textbf{(2) Monolithic ``generalist'' encoders do
not close this gap:} multi-domain contrastive training, physics-feature injection, and
channel-adversarial (DANN) objectives all fail a pre-registered bar against simple frozen
transfer, and we close this architecture line with the negative evidence. \textbf{(3) Native
per-protocol enrollment does close it where it matters:} on pooled multi-receiver BLE discovery, a
small natively trained encoder reaches 0.714 ARI where frozen transfer scores 0.284 and classical
features 0.193, and a variance decomposition shows why --- contrastive positives spanning
receivers drive collection (channel/receiver) variance to $\approx$0 while retaining 0.98 of
unit variance. These findings motivate a router-then-specialize system --- protocol
classification, per-protocol native encoders, and an explicit UNKNOWN tier --- which we implement
and validate end-to-end with regression-gated scenarios, including simultaneous WiFi + drone + BLE
discovery with zero cross-protocol identity collisions.
\end{abstract}

\begin{IEEEkeywords}
RF fingerprinting, open-set recognition, specific emitter identification, clustering,
metric learning, open-world discovery
\end{IEEEkeywords}

% =========================================================================
\section{Introduction}
% =========================================================================

Hardware imperfections --- oscillator offset, IQ imbalance, amplifier nonlinearity --- imprint a
device-specific signature on every RF emission. A decade of specific emitter identification (SEI)
work exploits this for \emph{closed-set} classification and, more recently, \emph{open-set
rejection}: flagging signals that match no enrolled device. But the operational question in
spectrum monitoring, counter-UAS, and interference hunting is usually different: \emph{how many
distinct emitters are out there, and which emissions belong to which?} The devices of interest are
precisely the ones that were never enrolled.

We formulate this as \emph{open-world emitter discovery}: an encoder $f$ maps raw IQ windows to
embeddings; at test time, emissions from a device population \emph{disjoint} from training are
clustered into identities, and clustering quality is scored against held-out true device labels.
Rejection is a by-product (an UNKNOWN route), not the goal.

This paper is deliberately an \emph{honest-evaluation} paper. Discovery metrics are fragile ---
they move with burst-integration length, clustering support conventions, receiver diversity, and
checkpoint selection --- and we found that several apparent effects in our own early results were
artifacts of exactly these knobs. All headline numbers therefore come from a locked protocol
(Sec.~\ref{sec:protocol}) with pre-registered success bands, device- and board-disjoint splits, a
one-touch rule for every held-out set, and negative results reported at the same standard as
positive ones.

\textbf{Contributions.}
\begin{enumerate}
  \item \textbf{A locked discovery-evaluation protocol} (Sec.~\ref{sec:protocol}): device-disjoint
  splits, a noise-scoring rule that stops density-based clusterers from being rewarded for
  refusing to cluster, separately reported \emph{assisted} (oracle-$K$) and \emph{deployable}
  ($K$-estimated) operating modes, and one-touch held-out sets.
  \item \textbf{The $N$-dependence of the cross-domain wall} (Sec.~\ref{sec:wall}): burst
  integration --- not representation --- carries most of cross-domain discovery. At $N{=}120$,
  frozen WiFi transfer, classical physics features, and native training converge to one band on
  8-way drone discovery; the representational gap that native training buys is
  $\approx$0.06--0.11 ARI.
  \item \textbf{Closure of the monolithic-generalist line} (Sec.~\ref{sec:generalist}): with
  pre-registered bands, multi-domain SupCon, physics-token injection, and channel-adversarial
  training all fail to beat frozen transfer at matched $N$; several actively hurt. We report the
  ablation ledger and close the line rather than tuning until something passes.
  \item \textbf{A three-way comparative study on same-model BLE discovery}
  (Sec.~\ref{sec:ble}): classical 19-D physics features vs.\ frozen WiFi transfer vs.\ a native
  1.49M-parameter encoder, on 31 same-model BLE units across 12 capture collections. The winner is
  \emph{regime-dependent}; in the pooled multi-receiver regime that deployment actually faces, the
  native encoder escapes the wall (0.714 vs.\ 0.284/0.193 ARI), and a variance decomposition
  identifies the mechanism (collection-invariance: collection $\eta^2\!\approx\!0$).
  \item \textbf{A router-then-specialize system} (Sec.~\ref{sec:system}): rate-aware protocol
  routing with a Mahalanobis UNKNOWN gate, per-protocol frozen native tiers of heterogeneous
  embedding dimension, per-protocol-group clustering with namespaced identities, validated by
  regression-gated scenarios including simultaneous three-protocol discovery with zero
  cross-protocol identity collisions.
\end{enumerate}

% =========================================================================
\section{Related work}
% =========================================================================
\label{sec:related}

\textbf{Open-set SEI (rejection).} OpenMax-style tail modeling, prototypical networks with
extreme-value calibration, and threshold-based verification reject signals from unenrolled
devices \todo{cite Bendale \& Boult 2016; Xie et al.\ 2021; Guo et al.\ 2024; survey refs from
LITERATURE\_SURVEY}. Discovery differs in kind: rejected signals must be further \emph{organized
into identities}, and the number of identities is itself unknown. We use rejection machinery only
as a routing gate (Sec.~\ref{sec:system}).

\textbf{Open-set discovery and novel-class clustering.} Generalized category discovery in vision
\todo{cite GCD/CiPR} clusters unlabeled classes at the \emph{category} level; recent UAV-signal
work \todo{cite arXiv 2603.24268} discovers signal \emph{types}. Our task is finer: distinct
\emph{physical units of the same model} --- same protocol, same PCB, same firmware --- separated
only by manufacturing tolerance.

\textbf{Same-model fingerprinting.} UAVSig \todo{cite MILCOM 2024} demonstrates closed-set
same-model drone identification; DRFF and BLE datasets \todo{cite dataset papers} provide
same-model populations we use in the open-world setting. Hiles et al.\ \todo{cite 2025 framework
paper} name emitter discovery as a task in a generic RFF framework; we instantiate and
stress-test it, including its failure modes.

\textbf{Cross-domain transfer in RFF.} Prior work reports both optimistic transfer and severe
receiver/day shift sensitivity \todo{cite WiSig paper; receiver-shift studies}. Our contribution
here is the \emph{controlled decomposition} of transfer into an integration effect and a
representation effect (Sec.~\ref{sec:wall}), which reconciles both readings.

% =========================================================================
\section{Problem formulation and locked evaluation protocol}
% =========================================================================
\label{sec:protocol}

\textbf{Task.} Train an encoder $f:\mathbb{R}^{2\times L}\!\to\!\mathbb{S}^{d-1}$ on labeled
windows from devices $\mathcal{D}_{\mathrm{train}}$. At evaluation, receive windows from
$\mathcal{D}_{\mathrm{test}}$, $\mathcal{D}_{\mathrm{test}}\cap\mathcal{D}_{\mathrm{train}}
=\emptyset$, grouped into \emph{bursts} (tracks) of $N$ windows by capture continuity --- never by
label. Embed, mean-pool each burst, L2-normalize, cluster. Score the partition against true device
labels, which are used \emph{only} for scoring.

\textbf{Metrics.} Adjusted Rand Index (ARI; chance $=0$) is primary; NMI and purity are
secondary. \textbf{Locked noise rule:} density-based clusterers may label points as noise
($-1$); every noise point is scored as its \emph{own singleton cluster}, and ARI/NMI/purity share
one support (all points). Without this rule, excluding noise from ARI/NMI rewards a clusterer for
refusing to cluster: under the old convention a margin-head variant scored 0.645--0.667 ARI while
discarding 45\% of bursts as noise; under the locked rule the same output scores 0.350. All
numbers in this paper use the locked rule.

\textbf{Assisted vs.\ deployable $K$.} We report two operating modes separately and never mix
them: \emph{assisted-$K$} (the operator supplies the emitter count; clustering is $K$-way ---
also the \emph{ceiling} diagnostic, reported for k-means and spectral clustering separately,
never max-of-methods) and \emph{deployable-$K$} (the system estimates $K$ by eigengap or HDBSCAN).
Count quality is reported as $K_{\mathrm{est}}$ and ARI at $K_{\mathrm{est}}$.

\textbf{Split discipline.} All splits are device-disjoint and, where fabrication structure is
known, \emph{board}-disjoint. Each domain has a development pool (all tuning) and a one-touch
held-out set evaluated exactly once with frozen hyperparameters and reported as-is. Success bands
for confirmatory runs are pre-registered before training. Splits, thresholds, and plan documents
are hash-frozen in the repository.

\textbf{Burst-construction audit.} Because bursts are built from storage-order continuity, we
verified they do not smuggle channel coherence into the score: \emph{scattered} bursts (windows
drawn from up to 72 distinct receiver$\times$day sessions of one device) score \emph{at least as
high} as coherent same-session bursts (e.g.\ margin-head 0.802 scattered vs.\ 0.396 coherent).
Burst-mean pooling therefore denoises toward a \emph{hardware} centroid rather than exploiting
within-session channel coherence.

% =========================================================================
\section{Datasets}
% =========================================================================
\label{sec:data}

\begin{table}[t]
\caption{Device populations. All evaluations are device-disjoint from all training.}
\label{tab:data}
\centering\small
\begin{tabular}{lllll}
\toprule
population & protocol & rate & devices & role \\
\midrule
WiSig ManyTx & WiFi 802.11a & 25\,MS/s & 150 (109/41) & train + in-domain \\
ORACLE & WiFi 802.11a & 5\,MS/s & 16 & aux.\ train (Sec.~\ref{sec:generalist}) \\
DRFF (mavicAir2) & DJI OcuSync & 50\,MS/s & 8 same-model & cross-domain eval \\
BLE D2 (XIAO) & BLE GFSK 1\,Mbps & 6\,MS/s & 31 same-model & Sec.~\ref{sec:ble} \\
\bottomrule
\end{tabular}
\end{table}

Table~\ref{tab:data} summarizes the populations. \textbf{WiSig}: 150 transmitters across 20
fabrication boards, 18 receivers $\times$ 4 capture days; board-disjoint 109 train / 41 discover
split, the 41 further partitioned into DEV (25 devices, 4 boards) and a one-touch TEST (16
devices, single board~18). \textbf{DRFF}: 8 DJI mavicAir2 airframes (same model) held out
entirely; 13 non-mavicAir2 airframes available for native training. \textbf{BLE D2}: 31 Seeed
XIAO ESP32-C3 units (same model), USRP B210 capture, 12 collections spanning wired/antenna links,
channels, distances, and two receivers; pre-registered split of 21 train / 2 val / 8 eval units
$\times$ disjoint collection sets. \textbf{ORACLE}: 16 USRP X310s at 5\,MS/s; used here only as a
training-pool ingredient, where it turns out to be informative by \emph{failing}
(Sec.~\ref{sec:generalist}).

% =========================================================================
\section{In-domain discovery and what makes it hard}
% =========================================================================
\label{sec:indomain}

\textbf{Encoders.} The WiFi specialist is a dual-branch encoder (time-domain and STFT branches,
fused; SupCon objective) producing 512-D embeddings; the same 1.49M-parameter architecture family
is used natively per domain throughout (drone: 4096-sample windows at 50\,MS/s; BLE: 1850-sample
windows at 6\,MS/s, 128-D head). No architecture search was performed per domain.

\textbf{Receiver diversity is the deployable lever.} On WiSig DEV, burst protocol dominates the
result. With bursts confined to one (receiver, day) session, discovery reaches 0.558 ARI
(embedding) / 0.368 (margin head); allowing bursts to span \emph{10 different receivers on the
same day} lifts the margin head to $0.772\pm0.037$ with $K_{\mathrm{est}}$ collapsing from
$\sim$43 to $16.3$ (true 18) and noise from 44\% to 0.8\%. Day diversity alone does \emph{not}
help (0.211). The per-receiver channel signature, not the day, is the confound; simultaneous
multi-sensor observation removes it. This is a deployable architecture requirement, not a data
trick --- and it shapes the system design in Sec.~\ref{sec:system}.

\textbf{Same-fabrication populations are the hard case.} The one-touch TEST (16 devices from a
single fabrication board) came back far below DEV: 0.219 vs.\ 0.727 ARI at the frozen operating
point, with every protocol row drifting 0.35--0.51. Local separability survives (kNN-1 up to
0.99) but the global 16-way partition of same-fab siblings collapses. Reported as-is: \emph{a
single fabrication batch of 16 same-model WiFi devices is roughly as hard to discover as
cross-domain drones}. Same-model discovery difficulty is fabrication-structure-dependent, which
is why the BLE study (Sec.~\ref{sec:ble}) uses 31 same-model units as its population.

% =========================================================================
\section{The cross-domain wall is $N$-dependent}
% =========================================================================
\label{sec:wall}

The canonical cross-domain cell: encoders that have never seen a drone must discover the 8
same-model mavicAir2 airframes. Table~\ref{tab:wall} shows the assisted-$K$ (oracle-$K$@8)
result as a function of burst integration $N$.

\begin{table}[t]
\caption{8-way drone discovery (ARI, assisted-$K$, k-means/spectral) vs.\ burst length $N$.
Fixed window cap 1500; the locked small-$N$ reference at cap 320 is 0.297 (see text).}
\label{tab:wall}
\centering\small
\begin{tabular}{lcc}
\toprule
method & $N{=}10$ (km) & $N{=}120$ (km / sp) \\
\midrule
frozen WiFi encoder (transfer) & 0.137 & 0.729 / 0.777 \\
classical 19-D physics features & 0.174 & 0.713 / 0.666 \\
native drone-trained encoder & 0.276 & \textbf{0.792 / 0.835} \\
multi-domain generalist (Sec.~\ref{sec:generalist}) & 0.197 & 0.625 / 0.666 \\
\bottomrule
\end{tabular}
\end{table}

Two readings. First, \textbf{burst integration washes out most of the transfer gap}: at
$N{=}120$, all four methods sit in one 0.62--0.79 band --- including hand-crafted features with
zero training. Heavy averaging toward the hardware centroid, not representation quality, carries
most of cross-domain discovery. Second, the \textbf{representational residual is real but
small}: native training buys $\approx$0.06--0.11 ARI over frozen transfer at matched $N$.

Small-$N$ numbers deserve a warning we learned the hard way: assisted-$K$ ARI at $N{=}10$ on
cross-domain features is \emph{burst-count-sensitive} --- the same frozen encoder scores 0.297 at
a 320-window cap but 0.137 at a 1500-window cap, purely because k-means over more noisy points
fragments more. We therefore treat $N{=}120$ as the cap-robust comparison and never quote a
small-$N$ cell as a fixed ``wall.''

Count estimation does \emph{not} come along for the transfer ride: at $N{=}120$ the native
stack recovers $K$ cleanly (eigengap $K{=}8$ on one seed; assisted success 1.0 in the demo
operating regime), while transferred features estimate $K{=}5$--7 with ARI@$K_{\mathrm{est}}$
$\le 0.63$. Good cluster geometry at oracle-$K$ does not imply a usable spectrum for $K$
estimation.

% =========================================================================
\section{Closing the generalist line}
% =========================================================================
\label{sec:generalist}

Could one protocol-agnostic encoder, trained across device populations with physics priors,
transfer better than frozen single-domain weights? We built the strongest version we could ---
121 training identities (WiSig 109 + ORACLE 12, zero drone data), a 1.84M-parameter encoder with
a 19-D physics-feature token and an LPC protocol-residual input channel, SupCon with
cross-condition positives, selection frozen before training --- and evaluated leave-one-protocol-out
on the 8 mavicAir2 airframes. The ledger:

\begin{itemize}
  \item \textbf{The generalist lost to everything at matched $N$} (Table~\ref{tab:wall}: 0.625
  at $N{=}120$ vs.\ frozen 0.729, native 0.792) and halved in-domain quality relative to the
  WiSig specialist (0.392 vs.\ $\approx$0.72) --- a cost that persisted under longer training and
  pool changes, i.e.\ intrinsic to the objective, not budget.
  \item \textbf{Physics-token injection hurt} both axes (removing it: transfer $+0.030$,
  in-domain $+0.122$), despite the token being demonstrably used during training (physics
  gradient share 0.06--0.33). Injected classical priors, fused by concat+FC, were not a free
  lunch.
  \item \textbf{The LPC residual channel was the one vindicated ingredient} (removing it:
  in-domain $-0.132$, transfer@$N{=}120$ $-0.136$) --- but it is rate-dependent: blind LPC
  extracts almost nothing at low oversampling (residual energy ratio 0.987 on 5\,MS/s OFDM vs.\
  0.001 at 50\,MS/s).
  \item \textbf{A degenerate domain poisons the pool.} ORACLE (12 identities, 5\,MS/s, 2
  conditions) never learned under the shared objective (its loss stayed flat while WiSig's fell)
  and removing it \emph{improved} transfer ($N{=}120$: $0.625\!\to\!0.781$). Data diversity must
  be individually learnable to help.
  \item \textbf{Channel-adversarial training (DANN), the strongest remaining lever, was
  pre-registered and returned NULL.} One run, bands frozen before training (SIGNAL: $N{=}10 \ge
  0.35$ or $N{=}120\ge0.80$): result 0.219 / 0.534 --- below both bands, and the $N{=}120$ number
  is the worst of every method measured. Weak adversarial pressure (adversary accuracy stayed
  0.63--0.94 against 0.25 chance) degraded transfer rather than enforcing useful invariance. No
  sweeps, no rescues, no second seed.
\end{itemize}

\textbf{Verdict.} With multi-domain SupCon, physics injection, residual views, positive-pair
placement, capacity, and adversarial invariance all exhausted against a frozen-transfer baseline
they never beat at matched $N$, we close the monolithic-generalist line on this pool. The two
levers left standing are genuine cross-family data diversity and \emph{per-domain native
enrollment behind a router} --- which the rest of the paper pursues. One positive finding
survives for future objectives: \emph{where} invariance pressure is placed is load-bearing ---
unconstrained positives gave the best transfer (0.806 at $N{=}120$) at the worst in-domain cost
(0.174) --- so the transfer/in-domain trade-off is a real dial, just not one any tested global
objective exploited productively.

% =========================================================================
\section{Same-model BLE discovery: a three-way comparative study}
% =========================================================================
\label{sec:ble}

If native enrollment is the answer, how much does it cost and what exactly does it buy? We ran a
pre-registered three-way comparison on BLE D2 (31 same-model units, 12 collections):
\textbf{A} --- 19 hand-crafted physics features (CFO, transient, modulation, spectral, IQ,
envelope families; zero training); \textbf{B1} --- the frozen WiSig WiFi encoder applied
cross-protocol (512-D); \textbf{B3} --- the same architecture trained natively from scratch on 21
train units (128-D; 4 runs, 35 min total on one GPU). A planned fine-tuning arm (B2) was skipped
by a pre-registered gate. Two evaluation regimes: \emph{canonical} (single capture collection ---
the easy, receiver-matched case) and \emph{pooled} (all eval collections mixed --- the regime a
multi-receiver base station actually faces).

\begin{table}[t]
\caption{BLE 8-unit discovery (ARI, assisted-$K$ unless noted), eval units never seen in
training. Pooled = collections mixed. $\mu$ over two seeds where shown.}
\label{tab:ble}
\centering\small
\begin{tabular}{lccc}
\toprule
cell & A (19-D) & B1 frozen (512-D) & B3 native (128-D) \\
\midrule
canonical, $N{=}1$ & 0.536 & 0.174 & \textbf{0.755} \\
canonical, $N{=}10$ & \textbf{0.925} & 0.597 & 0.892 \\
canonical, $N{=}120$ & 0.863 & \textbf{0.945} & 0.891 \\
pooled, $N{=}120$ & 0.193 & 0.284 & \textbf{0.714} \\
receiver transfer (fit R1, eval R2) & 0.571 & 0.454 & \textbf{0.759} \\
deployable HDBSCAN, canonical & 0.443 & \textbf{0.956} & 0.531 \\
eigengap $K_{\mathrm{est}}$ (true 8) & 3.2 & 11.0 & \textbf{8.6} \\
GPU ms/segment & --- & 0.638 & \textbf{0.142} \\
\bottomrule
\end{tabular}
\end{table}

\textbf{The winner is regime-dependent} (Table~\ref{tab:ble}). Classical features dominate
canonical low-integration ($N{=}10$: 0.925) --- they are per-segment discriminative. Frozen
transfer dominates canonical high-$N$ (0.945) and is the only method whose deployable HDBSCAN
recovers $K$ on canonical data (ARI 0.956) --- the textbook integration-backed transfer curve
from Sec.~\ref{sec:wall}, reproduced on a second protocol. \textbf{But in the pooled regime both
baselines hit the wall (0.193 / 0.284) and the native encoder escapes it: 0.714, $+$0.43 over
frozen transfer, at every $N$ (0.70 even at $N{=}1$), from as few as 5 training units.} On
canonical data at $N{=}120$ all three methods converge into a 0.082-wide band ---
representation choice barely matters on the easy regime and is decisive on the hard one.

\textbf{Mechanism.} A variance decomposition over embedding dimensions (multivariate $\eta^2$,
train units only) explains the pooled result:

\begin{table}[h]
\centering\small
\begin{tabular}{lccc}
\toprule
method & unit $\eta^2$ & collection $\eta^2$ & effect of per-collection calib. \\
\midrule
A (19-D) & 0.144 & 0.249 & rescues ($-0.225$ removed) \\
B1 (512-D) & 0.169 & 0.096 & helps ($-0.095$) \\
B3 (128-D) & \textbf{0.979} & \textbf{0.000} & none needed \\
\bottomrule
\end{tabular}
\end{table}

In raw physics features, receiver/channel (collection) variance \emph{exceeds} device variance
--- which is exactly why A collapses when collections are pooled and why per-collection
calibration rescues it. The native encoder's contrastive positives span collections inside every
batch, driving collection variance to $\approx$0 while unit variance is $\approx$all of it: it is
\emph{calibration-free by construction}, and it transfers across receivers best (0.759 vs.\
0.571/0.454). Neither learned encoder is a linear re-encoding of the physics features (ridge
$R^2 \le 0.12$ overall; the most recoverable family, carrier offset, caps at $R^2\!\approx\!0.3$)
--- the learned discriminative structure is largely orthogonal to the hand-crafted one.

\textbf{Negative micro-findings.} The BLE turn-on transient carries essentially no marginal
discovery signal on this dataset (dropping the transient family: $\Delta{=}-0.002$ ARI), a
dataset-limited refutation of a common narrowband-SEI intuition (the transient here is coarsely
sampled at 6\,MS/s and amplitude-normalized). And no method recovers $K{=}8$ by eigengap in the
pooled regime --- deployable count estimation remains open where it matters most.

% =========================================================================
\section{A router-then-specialize system}
% =========================================================================
\label{sec:system}

The evidence assembles into an architecture: \emph{protocol-classify first, then hand each track
to a native per-protocol encoder, cluster within protocol groups only, and route anything
unrecognized to an explicit UNKNOWN tier} --- the opposite of the monolithic generalist. We
implemented and validated this end-to-end.

\textbf{Pipeline.} (1) A tracker groups windows into per-emitter tracks from capture continuity
only. (2) A \emph{rate-aware} router classifies each track from 10 cheap spectral statistics
computed in absolute frequency (referenced to a common rate, so a 1\,MHz BLE burst reads as
narrowband whether captured at 6 or 25\,MS/s), with a per-class Mahalanobis gate (99th-percentile
training threshold) routing out-of-family signals to UNKNOWN. (3) A tier registry maps protocol
$\to$ frozen native encoder: drone (512-D), WiFi (512-D), BLE (128-D --- the B3 encoder of
Sec.~\ref{sec:ble}, shipped raw with no calibration, as EF-style variance analysis predicts), and
an UNKNOWN fallback tier. Mixed embedding dimensions never meet: clustering is always within a
protocol group, and identities are namespaced (\texttt{d1.ble}, \texttt{d2.drone}, \dots).
(4) Per track, $N^{*}{=}120$ windows are mean-pooled; per group, eigengap-$K$ (or operator-assisted
$K$) spectral clustering assigns identities.

\textbf{Validation.} Router held-out accuracy is 0.973 on 3-class training data and 1.000 on the
4-class held-out routing test (WiFi / drone / BLE / synthetic unknown; BLE side = eval units
only). The named adversarial risk --- a narrowband FM sweep now resembling narrowband BLE ---
was measured, not assumed: the sweep sits at Mahalanobis distance 78.6 from the BLE class (its
\emph{farthest}), 6.29 from the nearest class, above the 5.59 threshold, and routes to UNKNOWN.
Scenario gates (pre-set, 5 repeats each): single-protocol drone discovery at $K{=}3,4$;
mixed WiFi+drone; forced-UNKNOWN handling; BLE at $K\in\{2,3,4\}$ in both single-collection and
pooled-collection variants; and a grand-mixed scenario running 2 drone + 2 WiFi + 2 BLE emitters
simultaneously --- routing accuracy 1.00, per-protocol assisted discovery success 1.00, and
\emph{zero} cross-protocol identity collisions. Integration was regression-gated: after adding
the BLE tier and the rate-aware router, all pre-existing scenario results reproduced
byte-identically. Small-$K$ perfection here does not contradict the 8-unit battery numbers
(0.714 pooled): few same-model units at $N{=}120$ with a collection-invariant encoder is the
easy end of the curve, and we say so.

% =========================================================================
\section{Limitations}
% =========================================================================
\label{sec:limits}

(i) \textbf{Single dataset per claim-domain.} The BLE study is one dataset, one protocol; its
transient refutation and calibration hierarchy are dataset-limited. Zigbee/other narrowband PHYs
are enrolled-pending future work; until then such emitters route to UNKNOWN by design.
(ii) \textbf{Checkpoint selection was validation-starved} for the BLE native arm (2 val units;
selection ROC-AUC saturated from epoch 1, so early epochs won ties) --- the native encoder's
ceiling may be \emph{higher} than reported. (iii) \textbf{No temporal/aging axis}: BLE collections
are condition-disjoint, not day-disjoint; WiSig day diversity did not rescue discovery on its
own. (iv) \textbf{Deployable count estimation is unsolved in the pooled regime} (all methods
$\approx$0.19--0.26 deployable-$K$ ARI); assisted-$K$ is the shipping mode, honestly labeled.
(v) \textbf{Demo geometry is mocked} (single-receiver source data replayed round-robin);
localization/multilateration is out of scope. (vi) \textbf{Same-fab TEST hardness} (0.219 on
board-18) bounds expectations for dense single-batch populations; discovery claims should always
state the fabrication structure of the eval set.

% =========================================================================
\section{Conclusion}
% =========================================================================

Open-world emitter discovery is not open-set rejection with extra steps: its difficulty is
concentrated in regimes --- pooled receivers, same-fabrication populations, unknown $K$ --- that
rejection benchmarks never touch. Under a locked protocol we showed that burst integration, not
representation, carries most cross-domain discovery ($N$-dependent wall); that a carefully built
generalist encoder, physics injection, and adversarial invariance all fail pre-registered bars
against frozen transfer; and that cheap native per-protocol enrollment escapes the pooled-regime
wall ($0.284\!\to\!0.714$ on same-model BLE) by a measured mechanism --- learned
collection-invariance. The resulting router-then-specialize system is validated end-to-end with
regression gates and zero cross-protocol collisions. We offer the negative results, the locked
protocol, and the regime map as reusable equipment for the field: discovery numbers are only
meaningful with their integration length, $K$-mode, receiver structure, and fabrication
structure attached.

% =========================================================================
\section*{Reproducibility}
% =========================================================================
All splits, plan documents, and thresholds are hash-frozen; encoders are frozen artifacts with
recorded SHA-256; every table cell in this paper traces to a committed script and artifact path
in the project repository (FINDINGS.md, EXT\_FINDINGS.md, EVAL\_PROTOCOL.md ledgers).
\todo{public release decision + repo URL}

% =========================================================================
\begin{thebibliography}{9}
\bibitem{wisig} \todo{WiSig dataset citation}
\bibitem{oracle} \todo{ORACLE dataset citation}
\bibitem{drff} \todo{DRFF / drone RF dataset citation}
\bibitem{bled2} \todo{BLE XIAO dataset citation}
\bibitem{supcon} P.~Khosla et al., ``Supervised contrastive learning,'' NeurIPS 2020.
\bibitem{openmax} A.~Bendale, T.~Boult, ``Towards open set deep networks,'' CVPR 2016.
\bibitem{dann} Y.~Ganin et al., ``Domain-adversarial training of neural networks,'' JMLR 2016.
\bibitem{hdbscan} L.~McInnes, J.~Healy, ``Accelerated hierarchical density based clustering,'' ICDMW 2017.
\bibitem{uavsig} \todo{UAVSig MILCOM 2024 citation}
\end{thebibliography}

\end{document}
